\scanpage{58}
\mspage{57}
2 - 3 - \hfill n° 392

Montrons que c) implique d). Si c) est vraie, l'application $\underline a\longrightarrow\operatorname{div}\underline a$ est une bijection de $I(A)$ sur $D(A)$, (cf. § 1, n° 1) ; comme c'est un homomorphisme, (§ 1, n° 2), et que $D(A)$ est un groupe, tout élément de $I(A)$ est bien inversible.

Montrons que d) implique e). Si d) est vraie, tout idéal entier $\ne(0)$ de $A$ est de type fini, (chap. II, § 5, n° 6, th. 4), donc $A$ est noethérien ; comme $I(A)$ est un groupe, $D(A)$ est un groupe, et $A$ est donc complètement intégralement clos, (§ 1, n° 2, th. 1). Enfin, si $\underline p$ est un idéal premier non nul de $A$, et si $\underline m$ est un idéal maximal de $A$ contenant $\underline p$, l'anneau $A_{\underline m}$ est principal, (chap. II, § 5, n° 6, th. 4) ; comme $\underline pA_{\underline m}$ est premier non nul, on a nécessairement $\underline pA_{\underline m}=\underline mA_{\underline m}$, d'où $\underline p=\underline m$, (chap. II, § 2, n° 5, prop. 11), et $\underline p$ est maximal.

Montrons que e) implique f). En effet, si $\underline m$ est un idéal maximal de $A$, et si e) est vraie, $A_{\underline m}$ est un anneau noethérien intégralement clos, et son idéal maximal $\underline mA_{\underline m}$ est, ou bien $(0)$, ou bien le seul idéal premier non nul de $A_{\underline m}$ ; donc $A_{\underline m}$ est un corps ou un anneau de valuation discrète d'après la prop. 11 du § 1, n° 7.

Le fait que f) implique g) est évident.

Montrons enfin que g) implique a). Comme $A$ est l'intersection des $A_{\underline m}$, où $\underline m$ parcourt l'ensemble des idéaux maximaux, (chap. II, § 3, n° 3, cor. 4 du th. 1), g) implique que $A$ est intégralement clos et noethérien, donc que $A$ est un anneau de Krull, (§ 1, n° 3, cor. du th. 2). D'autre part, on montre que tout idéal premier non nul de $A$ est maximal comme dans la démonstration de d) $\Longrightarrow$ e). C. Q. F. D.

PROPOSITION 1.- \uline{Un anneau de Dedekind semi-local est principal}.

Soient $A$ un anneau de Dedekind semi-local, $K$ son corps des fractions, $p_1,\ldots,p_n$ ses idéaux maximaux, et $v_1,\ldots,v_n$ les valuations essentielles correspondantes ; ce sont les seules valuations essentielles de $A$. Soit $\underline a$ un idéal entier non nul de $A$. Puisqu'il est divisoriel, il existe, (§ 1, n° 4, prop. 5), des entiers $q_1,\ldots,q_n$ tels que $\underline a$ soit l'ensemble des $x\in K$ tels que $v_i(x)\ge q_i$ pour
